        \begin{tikzpicture}[shorten >= 1pt,
                auto,
                node distance=0.5cm,
                every node/.style = {shape=rectangle, draw, rounded corners},
                align = center]
            \node (q_linal)  {Некоторые понятия \\ линейной алгебры};

            \node (q_graphtheory) [below = of q_linal, ] {Некоторые сведения \\ из теории графов};
            \node (q_fortmallang) [right = 2 of q_graphtheory] {Общие сведения теории \\ формальных языков};

            \node (q_flpq) [below = of q_graphtheory] {Пути с ограничениями \\ в терминах формальных \\ языков};
            \node (q_reglang) [right = of q_flpq] {Регулярные \\  языки};
            \node (q_cflang) [right = of q_reglang] {Контекстно-свободные \\ языки и грамматики};
            \node (q_mcfl) [right = of q_cflang] {Многокомпонентные \\ контекстно-свободные \\ языки};

            \node (q_rpq) [below = 2 of q_flpq] {Поиск путей \\ с регулярными \\ ограничениями};
            \node (q_cfpq) [below = 2 of q_reglang] {Пути с ограничениями \\ в терминах  \\ контекстно-свободных \\ языков};
            \node (q_mcfpq) [right = of q_cfpq] {Пути с ограничениями \\ в терминах \\ многокомпонентных \\ контекстно-свободных \\ языков};

            \path[->]
            (q_linal)         edge (q_graphtheory)
            (q_graphtheory)   edge (q_flpq)
            (q_fortmallang)   edge (q_flpq)
            (q_fortmallang)   edge (q_reglang)
            (q_fortmallang)   edge (q_cflang)
            (q_fortmallang)   edge (q_mcfl)
            (q_reglang)       edge (q_rpq)
            (q_cflang)        edge (q_cfpq)
            (q_mcfl)          edge (q_mcfpq)
            (q_flpq)          edge (q_rpq)
            (q_flpq)          edge (q_cfpq)
            (q_flpq)          edge (q_mcfpq);
        \end{tikzpicture}
